%=============================================================================
% CASE STUDY 1 - COMP3308 (Introduction to Artificial Intelligence)
%-----------------------------------------------------------------------------
% Wired into the appendix block in main.tex after Appendix D, so it letters
% itself as Appendix E automatically via the custom \appendix patch.
%
% \ref keys below are wired to the real labels in the report:
%   sec:intro                       Sec I      Introduction            (unchanged)
%   subsec:system-architecture      Sec III-B  System Architecture
%   subsubsec:structured-retrieval  Sec III-C3 Structured Graph Retrieval
%   subsubsec:ir-extraction         Sec III-D2 IR Extraction (validation)
%   sec:data-results                Sec IV     Data and Results
%   sec:results-biencoder           Sec IV-C   Dense Bi-encoder Retrieval (label added there)
%   sec:analysis-discussion         Sec V      Analysis and Discussion
%   subsec:rq1 / subsec:rq2         Sec V-A / V-B   RQ1 / RQ2 analysis
%   subsec:limitations              Sec V-D    Limitations
%   sec:lr-dense / sec:lr-graph     Sec II-B1 / II-B2  Dense / Graph-augmented retrieval
%   sec:lr-nl2q / sec:lr-rag        Sec II-B3 / II-B4  NL-to-query / Query reformulation
%   app:prompts                     App C      Prompts and Tool Schemas (unchanged)
%   Tables: tab:kw-kpbiomed, tab:kw-pubmedake, tab:kw-stability,
%           tab:ir-reasoning, tab:ir-decomp, tab:ir-stratum, tab:ir-stability,
%           tab:xenc-results (full cross-encoder scores, in Appendix A)
%
% Citations resolve to existing .bib entries:
%   \cite{robertson1994bm25}   Robertson & Walker, SIGIR '94  (BM25)
%   \cite{haveliwala2002tspr}  Haveliwala, WWW 2002           (topic-sensitive PageRank)
%
% The Knuth & Moore alpha-beta ordering bound remains a commented block at the
% end of LO2; add a bib entry if the quantitative form of that claim is wanted.
%=============================================================================

\section{Case Study 1: COMP3308 (Introduction to Artificial Intelligence)}
\label{app:comp3308}

This case study evidences the learning outcomes of COMP3308, Introduction to
Artificial Intelligence, against the work documented in this report. In the
format agreed with the academic supervisor, an outcome the placement
exercised is evidenced by cross-reference to the body of the report, and new
material is written in place only where the placement did not itself
exercise the outcome, which is the case for LO2 alone. The mapping follows
the scoping recorded in the progress report.

\subsection{LO1: Problem formulation and search algorithms}
\textit{Formulate problem space description, select and apply suitable
search algorithms and analyse the issues involved.}

The problem-space formulation is the substance of Section~\ref{sec:intro},
which poses literature search as the reconciliation of two retrieval
paradigms, and of Section~\ref{subsubsec:structured-retrieval}, which gives the structured
stage a formal statement in which a query is a typed entity list with a
recursive boolean expression over it and candidacy is defined by the
satisfaction relation $w \models \varphi$. Search-algorithm selection and
application run through the cascade of
Section~\ref{subsec:system-architecture}, namely BM25 over per-entity-type
text indexes for non-exact entity matching, direct property match for exact
entities, top-100 dense retrieval by cosine similarity, and cross-encoder
re-ranking. The issues analysed alongside include the restricted-universe
semantics required for NOT (Section~\ref{subsubsec:structured-retrieval}),
the recall bound the boolean gate places on first-stage retrieval
(Section~\ref{subsec:limitations}), and the revision of the formulation
itself, from the fuzzy-match-and-traverse design recorded in the progress
report to the compiled subquery-per-entity design documented here.

\subsection{LO2: Minimax search and alpha-beta pruning}
\textit{Understand and apply minimax search and alpha-beta pruning in game
playing.}

The placement contained no adversarial search, so no minimax or alpha-beta
implementation forms part of the system and none is claimed here. The
progress report scoped the outcome accordingly, and this subsection
addresses it at the level of the principle those algorithms share with the
deployed architecture.

That principle is staged elimination ordered by cost. Alpha-beta's
contribution to minimax is to discard subtrees that provably cannot
influence the root decision, so that search effort concentrates where it
can still change the outcome, and the retrieval cascade of
Section~\ref{subsec:system-architecture} is built on the same economy. Its first cut,
the boolean gate, discards every work failing the query expression, and
this pruning is exact in the same sense as an alpha-beta cutoff. With
respect to the objective the system defines, a work failing $w \models E$
is by construction not an answer, so no later stage could have surfaced it.
The analogy must, however, be stated with its limits. Alpha-beta is
lossless with respect to the minimax value, whereas the cascade's second
cut, the truncation of the ranked pool to the top 100 candidates, is not,
because a work ranked just outside the cut by the bi-encoder might have
been ranked first by the cross-encoder and nothing bounds that loss. The
truncation is therefore heuristic forward pruning, closer in character to
beam search than to alpha-beta, and the gate's own exactness holds against
the system's definition of candidacy rather than an external guarantee of
relevance, the caveat Section~\ref{subsec:limitations} records.

The comparison also carries over alpha-beta's best-known practical
property. The cutoffs alpha-beta achieves are governed by move ordering, in
that the earlier a strong move is examined, the more of the tree the same
bounds eliminate. The cascade's analogue is the ranking fidelity of its
cheap stage. The better the bi-encoder orders the candidate pool, the less
relevance the fixed truncation discards, which is the property the
bi-encoder benchmarking of Section~\ref{sec:results-biencoder} measures and
the selection of Section~\ref{subsec:rq2} optimises for. The outcome is
therefore evidenced through understanding, of cost-ordered elimination, of
the distinction between sound and heuristic pruning, and of the dependence
of pruning safety on ordering quality, rather than through a game-playing
implementation.
% Optional strengthening for the ordering point: under perfect move
% ordering alpha-beta examines O(b^{d/2}) leaves against minimax's O(b^d)
% (Knuth & Moore, "An analysis of alpha-beta pruning", Artificial
% Intelligence 6(4), 1975). Add the bib entry if you want the quantitative
% form of the claim.

\subsection{LO3: Supervised, unsupervised, and probabilistic reasoning}
\textit{Understand the basic principles and analyse the strengths,
weaknesses and applicability of some of the main AI algorithms for
supervised learning, unsupervised learning and probabilistic reasoning.}

The keyword benchmark compares the supervised and unsupervised families
directly, running supervised fine-tuned models (KeyBART, KBIR-inspec)
against unsupervised statistical and graph-based extractors (TF-IDF,
YAKE, TextRank, MultipartiteRank), with the deployed prompted language
model evaluated alongside both
(Tables~\ref{tab:kw-kpbiomed}--\ref{tab:kw-pubmedake}). The analysis in
Section~\ref{subsec:rq1} then diagnoses the characteristic weakness of the
supervised family when KeyBART's lead on its own training distribution,
OAGKX, collapses off it, a case of overfitting read from cross-dataset
behaviour. Section~\ref{subsec:rq2} extends the same analysis to retrieval,
where the citation-supervised SPECTER2 variants top their in-domain task
yet finish near the bottom on average, and the biomedically supervised
MedCPT fails to lead even on biomedical text while general-purpose models
hold up across all tasks, an applicability analysis of specialisation
against generality.

Probabilistic reasoning enters at the system's core and in its evaluation.
The matcher of the structured stage is BM25 (Section~\ref{subsubsec:structured-retrieval}),
grounded in the probabilistic relevance framework~\cite{robertson1994bm25}.
Section~\ref{sec:lr-graph} analyses topic-sensitive
PageRank~\cite{haveliwala2002tspr}, a random-walk model whose teleport
distribution conditions relevance propagation on the query, as the
strongest conceptual fit among graph methods. Every model selection in
Section~\ref{sec:data-results} is then made under bootstrap uncertainty
estimates, with overlapping intervals read as ties and the stochasticity of
temperature-zero endpoints quantified over repeated runs
(Tables~\ref{tab:kw-stability} and~\ref{tab:ir-stability}).

\subsection{LO4: Designing, implementing and evaluating AI algorithms}
\textit{Gain practical experience in designing, implementing and evaluating
AI algorithms.}

The design and implementation experience is concentrated in the three
artefacts of the structured stage: firstly, the natural-language-to-IR
extractor, with its system prompt, forced tool-call schema and validation
step (Section~\ref{subsubsec:structured-retrieval},
Appendix~\ref{app:prompts}); secondly, the keyword-generation module, whose
prompt and configuration Appendix~\ref{app:prompts} reproduces; and
thirdly, the deterministic compiler that turns the validated representation
into Cypher. The evaluation experience spans building instruments as well
as using them. A 116-query stratified gold set was constructed by hand for
a task no existing benchmark covers (Section~\ref{subsubsec:ir-extraction}),
and the component benchmarks report bootstrap confidence intervals
throughout (Tables~\ref{tab:kw-kpbiomed}--\ref{tab:xenc-results}), a
reasoning ablation (Table~\ref{tab:ir-reasoning}), per-component and
per-stratum decompositions (Tables~\ref{tab:ir-decomp}
and~\ref{tab:ir-stratum}), and run-to-run stability checks for the
non-deterministic components (Tables~\ref{tab:kw-stability}
and~\ref{tab:ir-stability}). These evaluations carried deployment stakes,
since the models they selected are the ones serving the industry partner's
production system.

\subsection{LO5: Presenting and interpreting data in verbal and written form}
\textit{Present and interpret data and information in verbal and written
form.}

The written component comprises the progress report and this report, whose
Sections~\ref{sec:data-results} and~\ref{sec:analysis-discussion} present the benchmark
data and interpret it against the research questions, including the
disciplines of reading overlapping confidence intervals as ties and of
decomposing a headline accuracy into its binding component. The verbal
component has run continuously through the placement, with benchmark
results and design decisions presented at the industry partner's weekly
stand-ups and in regular academic-supervisor meetings as they were
produced. The concluding verbal presentation, the ESIPS oral to the
academic supervisor and the industry partner, is scheduled for
14 August 2026.

\subsection{LO6: Achievements and shortcomings of AI, and links to other areas}
\textit{Appreciate some of the main ideas and views in AI, achievements and
shortcomings of AI and the links between AI and other computer science
areas.}

The literature review is a sustained critical appraisal of modern AI's
record on retrieval. Section~\ref{sec:lr-dense} weighs dense retrieval's
achievements against the share of relevant papers it still misses on
realistic scientific queries, Section~\ref{sec:lr-rag} presents the
evidence that LLM query expansion degrades precisely the strong retrievers
it would need to help, and Section~\ref{sec:lr-graph} separates conceptual
promise from empirical record for graph methods. The same appraisal is
turned on this system's own components in Section~\ref{subsec:limitations},
which records the costs of building artefacts with language models, namely
occasional generation of unsupported keyphrases, imperfect run-to-run
reproducibility, and dependence on hosted endpoints that can change without
notice. The links between AI and neighbouring computer-science areas are
load-bearing in the architecture rather than incidental. Intent extraction
meets database query compilation in the IR-to-Cypher design
(Sections~\ref{sec:lr-nl2q} and~\ref{subsubsec:structured-retrieval}), graph structure and
information retrieval meet in the typed knowledge graph, and
Section~\ref{subsec:rq2} turns model selection into a systems-engineering
question of memory footprint, latency and licensing.